\documentclass[11pt]{article}
\usepackage[utf8]{inputenc}
\usepackage[T1]{fontenc}
\usepackage{lmodern}
\usepackage[a4paper,margin=1in]{geometry}
\usepackage{booktabs}
\usepackage{longtable}
\usepackage{hyperref}
\usepackage{graphicx}
\usepackage{enumitem}

\title{Building a Reproducible Katharevousa Greek NLP Pipeline:\\
From OCR Reconstruction and LLM-Assisted Gold Data to Dependency Parsing}
\author{George Mikros \and Fotios Fitsilis}
\date{\today}

\begin{document}
\maketitle

\begin{abstract}
Katharevousa Greek remains under-served by modern NLP tools despite its central role in archival and parliamentary corpora. We present a reproducible end-to-end pipeline for building a dependency parsing benchmark and model stack for Katharevousa from noisy historical sources. Starting from OCR-derived parliamentary material, we apply targeted text reconstruction, large-scale LLM-assisted annotation, and quality-controlled CoNLL-U consolidation to produce a frozen gold dataset of 1,697 sentences. We evaluate feature-based, transformer, mBERT, and Stanza variants on a fixed held-out split. The best overall model is an XLM-R parser with UPOS 0.8893, DEPREL F1 0.7250, UAS 0.6098, and LAS 0.5162. We release scripts and evaluation artifacts to support reproducibility.
\end{abstract}

\section{Introduction}
Historical and formal Greek registers, especially Katharevousa, pose a structural mismatch for modern Greek NLP systems trained primarily on Demotic corpora. This affects morphology and syntax, especially dependency structure in long institutional prose. We document a reproducible research workflow from OCR reconstruction to parser selection and release readiness.

\section{Data and Pipeline}
The pipeline builds on OCR-derived post-junta parliamentary material and includes:
\begin{itemize}[leftmargin=1.25em]
  \item reconstruction-aware preprocessing (hyphenation and sentence restoration),
  \item schema-constrained LLM annotation with validation and retries,
  \item deterministic freeze of final CoNLL-U and manifests,
  \item fixed train/test split evaluation with shared metrics.
\end{itemize}

Frozen snapshot summary: 1,697 sentences total, with 1,565 from first-pass batches and 132 from retry replacements.

\section{Experimental Setup}
Fixed split (seed 42): train 1,357 sentences, test 340 sentences (4,093 tokens). Metrics: UPOS accuracy, DEPREL weighted F1, UAS, LAS.

\subsection{Compared Models}
\begin{itemize}[leftmargin=1.25em]
  \item Logistic baseline (feature-based parser).
  \item Transformer XLM-R (v1\_opt).
  \item mBERT variants.
  \item Stanza pretrained Greek baseline.
  \item Stanza custom-trained baseline.
\end{itemize}

\section{Results}
\begin{table}[h]
\centering
\caption{Main benchmark results on fixed test split.}
\begin{tabular}{lcccc}
\toprule
Model & UPOS & DEPREL F1 & UAS & LAS \\
\midrule
Logistic baseline (v3) & 0.9040 & 0.7451 & 0.5781 & 0.5072 \\
XLM-R (v1\_opt) & 0.8893 & 0.7250 & \textbf{0.6098} & \textbf{0.5162} \\
mBERT (best) & 0.8260 & 0.6076 & 0.5886 & 0.4537 \\
Stanza pretrained & 0.6125 & 0.4242 & 0.6079 & 0.3396 \\
Stanza custom (600 steps) & 0.7694 & 0.6588 & 0.5756 & 0.4943 \\
\bottomrule
\end{tabular}
\end{table}

XLM-R offers the best structural parsing performance (UAS/LAS), while the logistic baseline remains competitive in UPOS/DEPREL at this dataset scale. Custom Stanza improves substantially over pretrained Stanza but remains below XLM-R overall.

\section{Discussion}
The experimentation journey highlights practical lessons: sensitivity to scheduler/weight tuning, fragile gains from hard-case augmentation, and framework-level constraints in strict UD handling. Negative results are reported because they inform reproducible low-resource methodology.

\section{Conclusion}
The pipeline demonstrates that release-ready Katharevousa parsing is feasible with constrained annotation budgets when LLM-assisted generation is paired with strict validation and reproducible evaluation. The XLM-R v1\_opt model is the recommended first public release.

\bibliographystyle{plain}
\bibliography{references}

\end{document}
